\lecture{6}{21. Oktober 2025}{Fri svingning/egensvingning af MDOF-systemer (del I)}

\begin{problemwithimage}{p6_1}{6.1}
  Vi betragter det viste masse-fjeder-system, der skal analyseres som et system med to frihedsgrader; navnligt $d_1(t)$ og $d_2(t)$.
  a) Opstil masse- og stivhedsmatricen for systemet.
  b) Bevis, at systemets udæmpede vinkelegenfrekvenser er givet ved
  \[ 
  \omega_1 = \sqrt{\frac{k}{m}}, \qquad \omega_2 = \sqrt{5}\omega_1
  .\]
  c) Bestem de tilhørende egensvingningsformer/-vektorer, $\phi_1$ og $\phi_2$.
\end{problemwithimage}

\begin{figure}[ht]
    \centering
    \incfig[0.8]{p6_1-1}
    \caption{Frit-legeme diagram for de to masser i systemet.}
    \label{fig:p6_1-1}
\end{figure}

\begin{derivation}
  På \Cref{fig:p6_1-1} ses et frit-legeme diagram for hver af de to masser i systemet.
  På baggrund af dette kan Newtons 2. lov benyttes for først den første masse, som
  \begin{align*}
    m \ddot{d}_1(t) &= - kd_1 + 2 k \left( d_2 - d_1 \right) \\
                    &= 2k d_2 - 3k d_1 \\
    \implies \quad m \ddot{d}_1 (t) + 3k d_1 - 2k d_2 &= 0
  .\end{align*}
  Og dernæst for den anden masse, som
  \begin{align*}
    m \ddot{d}_2(t) &= - 2k \left( d_2 - d_1 \right) - k d_2 \\
          &= 2k d_1 - 3k d_2 \\
    \implies m \ddot{d}_2 (t) + 3k d_2 - 2k d_1 &= 0
  .\end{align*}
  Dermed har vi de to ligninger
  \begin{align*}
    m \ddot{d}_1 (t) + 3k d_1 - 2k d_2 &= 0 \\
    m \ddot{d}_2 (t) + 3k d_2 - 2k d_1 &= 0
  .\end{align*}
  Dette giver systemet
  \[ 
  \begin{bmatrix}
  m & 0\\
  0 & m\\
  \end{bmatrix} \begin{bmatrix}
  \ddot{d}_1 (t)\\
  \ddot{d}_2 (t)\\
  \end{bmatrix} + \begin{bmatrix}
  3k & - 2k\\
  -2k  & 3k \\
  \end{bmatrix} \begin{bmatrix}
  d_1\\
  d_2\\
  \end{bmatrix} = \begin{bmatrix}
  0\\
  0\\
  \end{bmatrix}
  .\]
  Heraf ses at masse- og stivhedsmatricerne er hhv.
  \begin{align*}
    M &= \begin{bmatrix}
    m & 0\\
    0 & m\\
    \end{bmatrix} \\
      K &= \begin{bmatrix}
      3k & -2k\\
      -2k & 3k\\
      \end{bmatrix}
  .\end{align*}
  Vi har at den udæmpede vinkelegenfrekvens $\omega$ kan findes ved
  \[ 
  \mathrm{det} \left( K - \omega^2 M \right) = 0
  .\]
  Ved indsættelse fås
  \begin{equation*}
    \mathrm{det}\left( \begin{bmatrix}
    3k - \omega^2 m &  - 2k \\
    - 2k & 3k - \omega^2 m\\
    \end{bmatrix} \right) = \left( 3k - \omega^2 m \right)^2 - 4 k^2 = 0
  .\end{equation*}
  De positive løsninger for $\omega$ af denne er hhv.
  \[ 
    \omega_1 = \sqrt{\frac{k}{m}}, \qquad \omega_2 = \sqrt{5} \cdot \sqrt{\frac{k}{m}} = \sqrt{5} \omega_1
  .\]
  Altså er det vist.

  For hver af disse 2 udæmpede vinkelegenfrekvenser findes en egenvektor $\phi_j$ givet ved
  \[ 
  \left( K - \omega_j^2 M \right) \phi_j = 0
  .\]
  For $\omega_1 = \sqrt{k / m}$ fås
  \begin{align*}
    0 &= \begin{bmatrix}
    3k - \frac{k}{m}m & - 2k\\
    -2k & 3k - \frac{k}{m}m\\
    \end{bmatrix} \phi_1 \\
      &= \begin{bmatrix}
    2k & -2k\\
    -2k & 2k\\
    \end{bmatrix} \begin{bmatrix}
    \phi_{11}\\
    \phi_{21}\\
    \end{bmatrix} \\
      &= 2k \left( \phi_{11}- \phi_{21} \right) \\
      \implies \qquad \phi_{11} &= \phi_{21}
  .\end{align*}
  Hvis vi lader $\phi_{11} = 1$ så $\phi_1 = \begin{bmatrix}
  1\\
  1\\
  \end{bmatrix}$.
  For $\omega_2 = \sqrt{5}\omega_1$ får vi
  \begin{align*}
    0 &= \begin{bmatrix}
    3k - 5 \frac{k}{m}m & - 2k\\
    -2k & 3k - 5 \frac{k}{m}m\\
    \end{bmatrix}\phi_2  \\
      &= \begin{bmatrix}
    -2k & -2k\\
    -2k & -2k\\
    \end{bmatrix} \begin{bmatrix}
    \phi_{12}\\
    \phi_{22}\\
    \end{bmatrix} \\
      &= - 2k \left( \phi_{12} + \phi_{22} \right) \\
      \implies \qquad \phi_{12} &= - \phi_{22}
  .\end{align*}
  Vi lader $\phi_{12} = 1$ så $\phi_2 = \begin{bmatrix}
  1\\
  -1\\
  \end{bmatrix}$
\end{derivation}

\finalresult{
  \begin{aligned}
    \mathrm{a}) &\quad M = \begin{bmatrix}
    m & 0\\
    0 & m\\
    \end{bmatrix}, \quad K = \begin{bmatrix}
    3k & -2k\\
    -2k & 3k\\
    \end{bmatrix} \\
    \mathrm{b}) &\quad \omega_1 = \sqrt{\frac{k}{m}}, \quad \omega_2 = \sqrt{5}\omega_1 \\
    \mathrm{c}) &\quad \phi_1 = \begin{bmatrix}
    1\\
    1\\
    \end{bmatrix}, \quad \phi_2 = \begin{bmatrix}
    1\\
    -1\\
    \end{bmatrix}
  \end{aligned}
}

\begin{problemwithimage}{p6_2}{6.2}
  Vi betragter det viste rammesystem, hvori de horisontale/tværgående bjælker er uendeligt stive og hver har massen $m = \qty{1,5}{\kg}$, mens de vertikale bjælker er masseløse og hver har længden $L = \sqrt[3]{\num{1,2}} \, \unit{\m}$ og stivheden $k = \frac{12 EI}{L^3}$ \, \unit{\N\per\m} med $EI = \qty{10}{\N\m\squared}$.
  Systemet skal modelleres som et system med to frihedsgrader, hvor den første frihedsgrad, $d_1(t)$, er den horisontale flytning af den øverste tværgående bjælke, mens den anden frihedsgrad, $d_2(t)$, er den horisontale flytning af den nederste tværgående bjælke.
  Begyndelsesbetingelserne $d_0 = \phi_2$ og $\dot{d}_0 = 0$ påføres, hvortil det bemærkes, at $\phi_2$ er den anden egensvingningsform.
  a) Opstil stivheds- og massematricen for systemet.
  b) Bestem systemets udæmpede vinkelegenfrekvenser, $\omega_j$, og de tilhørende egensvingningsformer, $\phi_j$.
  c) Bestem $d_1(t)$ og $d_2(t)$.
\end{problemwithimage}

\begin{figure}[ht]
    \centering
    \incfig[0.85]{p6_2fbd}
    \caption{Frit-legeme diagram for systemet}
    \label{fig:p6_2fbd}
\end{figure}

\begin{derivation}
  Systemet kan betrates som vist på \Cref{fig:p6_2fbd}.
  Benyttes Newtons 2. lov for den første masse fås:
  \begin{align*}
    m \ddot{d}_1(t) &= - 2k (d_1 - d_2) \\
    \implies \quad m \ddot{d}_1 (t) + 2k d_1 - 2k d_2 &= 0
  .\end{align*}
  Gøres det samme for den anden masse fås:
  \begin{align*}
    m \ddot{d}_2(t) &= 2k \left( d_1 - d_2 \right) - 2k d_2 \\
    \implies \quad m \ddot{d}_2 (t) + 4 d_2 k - 2 d_1 k &= 0
  .\end{align*}
  Dermed har vi ligningerne:
  \begin{align*}
    m \ddot{d}_1 (t) + 2k d_1 -  2k d_2 &= 0 \\
    m \ddot{d}_2(t) + 4 d_2 k - 2 d_1 k &= 0
  .\end{align*}
  Hvilket giver systemet:
  \[ 
  \begin{bmatrix}
  m & 0\\
  0 & m\\
  \end{bmatrix} \begin{bmatrix}
  \ddot{d}_1 (t)\\
  \ddot{d}_2(t)\\
  \end{bmatrix} + \begin{bmatrix}
  2k & - 2k\\
  - 2k & 4k \\
  \end{bmatrix} \begin{bmatrix}
  d_1(t)\\
  d_2(t)\\
  \end{bmatrix} = \begin{bmatrix}
  0\\
  0\\
  \end{bmatrix}
  .\]
  Dermed er masse- og stivhedsmatricerne hhv:
  \begin{gather*}
    M = \begin{bmatrix}
    m & 0\\
    0 & m\\
    \end{bmatrix} = \begin{bmatrix}
    \num{1,5} & 0\\
    0 & \num{1,5} \\
    \end{bmatrix} \, \unit{\kg}\\
      K = \begin{bmatrix}
      2k & -2k\\
      -2k & 4k\\
      \end{bmatrix} = 2k \begin{bmatrix}
      1 & -1\\
      -1 & 2\\
      \end{bmatrix} = \frac{24 EI}{L^3} \begin{bmatrix}
      1 & -1\\
      -1 & 2\\
      \end{bmatrix} = \begin{bmatrix}
      200 & -200\\
      -200 & 400\\
      \end{bmatrix}\, \unit{\N\per\m}
    .\end{gather*}
    Den udæmpede vinkelegenfrekvens $\omega$ kan findes ved
    \[ 
    \mathrm{det} \left( K - \omega^2 M \right) = 0
    .\]
    Ved indsættelse fås
    \begin{align*}
      0 =& \mathrm{det} \left( \begin{bmatrix}
    \qty{200}{\N\per\m} - \omega^2 \cdot \qty{1,5}{\kg} & \qty{-200}{\N\per\m} \\
    \qty{-200}{\N\per\m} & \qty{400}{\N\per\m} - \omega^2 \cdot \qty{1,5}{\kg} \\
\end{bmatrix} \right) \\
          =& \left( \qty{200}{\N\per\m} - \omega^2 \cdot \qty{1,5}{\kg} \right) \left( \qty{400}{\N\per\m} - \omega^2 \cdot \qty{1,5}{\kg}  \right) - \left( \qty{-200}{\N\per\m}  \right) \left( \qty{-200}{\N\per\m}  \right) \\
          \implies& \qquad \omega_1 = \qty{7,14}{\s^{-1}} \quad \omega_2 = \qty{18,68}{\s^{-1}}
    .\end{align*}
    Her kan vi nu løse
    \[ 
      \left( K - \omega^2_j M \right)\phi_j = 0
    .\]
    Vi får hhv.
    \begin{align*}
      \phi_1 &= \begin{bmatrix}
      - \num{0,618}\\
      1\\
      \end{bmatrix} \\
        \phi_2 &= \begin{bmatrix}
        \num{1,618} \\
        1\\
        \end{bmatrix}
    .\end{align*}
    Vi har den generelle løsningsformel:
    \[ 
    d(t) = \phi_1 \left( A_1^{(1)} \cos \left( \omega_1 t \right) + A_1^{(2)} \sin \left( \omega_1 t \right) \right) + \phi_2 \left( A_2^{(1)} \cos \left( \omega_2 t \right) + A_2^{(2)} \sin \left( \omega_2 t \right) \right)
    .\]
    Begyndelsesbetingelsen $d_0 = \phi_2$ giver os
    \begin{align*}
      \phi_2 =& \phi_1 A_1^{(1)} + \phi_2 A_2^{(1)} \\
      \implies& \quad A_1^{(1)} = 0, \quad A_2^{(1)} = 1
    .\end{align*}
    Dermed har vi
    \[ 
    d(t) = \phi_1 A_1^{(2)} \sin \left( \omega_1 t \right) + \phi_2 \left( \cos \left( \omega_2 t \right) + A_2^{(2)} \sin \left( \omega_2 t \right) \right)
    .\]
    Vi differentierer dette for at få
    \[ 
    \dot{d} (t) = \omega_1 \phi_1 A_1^{(2)} \cos \left( \omega_1 t \right) + \phi_2 \left( - \omega^2 \sin \left( \omega_2 t \right) + \omega_2 A_2^{(2)} \cos \left( \omega_2 t \right) \right)
    .\]
    Ved indsættelse af vores begyndelsesbetingelse, $\dot{d}_0 = 0$ fås
    \begin{align*}
      0 =& \omega_1 \phi_1 A_1^{(2)} + \phi_2 \omega_2 A_2^{(2)} \\
      \implies& \quad A_1^{(2)} = A_2^{(2)} = 0
    .\end{align*}
    Dermed fås:
    \[ 
    d(t) = \phi_2 \cos \left( \omega_2 t \right)
    .\]
    Hvilket svarer til systemet
    \[
      d(t) = \begin{bmatrix}
      \num{1,62}\\
      1\\
      \end{bmatrix} \cos \left( \qty{18,7}{\s^{-1}} \cdot t \right)
    .\]
  \end{derivation}

\finalresult{
\begin{aligned}
    \mathrm{a}) &\quad M = \begin{bmatrix}
    \num{1,5}  & 0\\
    0 & \num{1,5} \\
    \end{bmatrix} \, \unit{\kg}, \quad  K = \begin{bmatrix}
    200 & -200\\
    -200 & 400\\
    \end{bmatrix} \, \unit{\N\per\m} \\
    \mathrm{b}) &\quad \omega_1 = \qty{7,14}{\s^{-1}}, \quad  \omega_2 = \qty{18,7}{\s^{-1}}, \quad  \phi_1 = \begin{bmatrix}
    \num{-0,618} \\
    1\\
    \end{bmatrix}, \quad \phi_2 = \begin{bmatrix}
    \num{1,62} \\
    0\\
    \end{bmatrix} \\
    \mathrm{c}) &\quad d(t) = \begin{bmatrix}
    \num{1,618} \\
    1\\
    \end{bmatrix} \cos \left( \qty{18,7}{\s^{-1}} \cdot t \right)
.\end{aligned}
}
